\begin{lemma}{Numeriklemma}{Numeriklemma}
Seien 
\[h_1: [1, \infty[ \rightarrow \R, 
t\mapsto\frac{\sqrt t}{t+1}
\bigg((\pi-1)\sqrt{\log(t)}+\sqrt{\log(t)+2\pi}\bigg)\text,\]
\[h_2: \R_{>0} \rightarrow \R,
t \mapsto \frac{\sqrt{t+1}}t
\bigg( (\pi-1)\sqrt{\log(1+t)} + \sqrt{\log(1+t) + 2\pi} \bigg)\text.\]
Dann gilt:
\begin{enumerate}[(1)]
	\item \label{Numeriklemma0}
	$\forall\, t \in \R_{>0}: \log(t) = 2\Sum{k=0}\infty \frac 1 {2k+1}
	\left(\frac{t-1}{t+1} \right)^{2k+1}$.
	\item \label{Numeriklemma1}
	$h_1 < \sqrt{2 \pi}$.
	\item \label{Numeriklemma2}
	$h_2$ ist monoton fallend und für alle $t \in [\frac{13}2, \infty[:
	h_2(t) < \sqrt{2 \pi}$.
	\item\label{Numeriklemma4}
	$\forall\, \eta\in\!\  ]0,1[ 
	 \forall \,t\in 
	 \bigg[\frac 4\eta \log\left(\frac 2 \eta\right), \infty\bigg[
	 : \log(2t)< \eta t$.
\end{enumerate}
\end{lemma}
\begin{proof}
(1) Sei $t \in \R_{>0}$.\\
Die Taylorentwicklung des Logarithmus' liefert
\begin{align}\forall\, x \in\!\ ]-1, 1[: \log(1+x) = \sum_{k=0}^\infty (-1)^k 
\frac {x^{k+1}}{k+1}\text. \label{nleq:1}\end{align}
Sei $x_t \coloneqq  \frac{t-1}{t+1}$. Dann gilt:
\[ \frac{1+x_t}{1-x_t} = \frac{1+\frac{t-1}{t+1}}{1-\frac{t-1}{t+1}}
= \frac{\frac{t+1+t-1} {t+1}}{\frac{t+1-t+1}{t+1}}
= \frac{\frac{2t}{t+1}}{\frac 2 {t+1}}=t \text.\]
Ferner gilt: $-t < t$, sodass $-(t+1)=-t -1< t-1$, also $x_t > -1$,
sowie $t-1 < t+1$, sodass $x_t < 1$. Also $x_t, -x_t \in\!\  ]-1,1[$.\\
Erhalte:
\begin{align*}
\log(t) &= \log\left(\frac{1+x_t}{1-x_t} \right) = \log(1+x_t)-\log(1-x_t)\\
&\overset{\eqref{nleq:1}}=
\left(\sum_{k=0}^\infty (-1)^k\frac{x_t^{k+1}}{k+1} \right)
-\left(\sum_{k=0}^\infty (-1)^k(-1)^{k+1}\frac{x_t^{k+1}}{k+1} \right)\\
&=\sum_{k=0}^\infty\left((-1)^k-(-1)^k(-1)^{k+1} \right)\frac{x_t^{k+1}}{k+1}\\
&=2\sum_{\substack{k=0,\\ k \text{ gerade}}}^\infty\frac{x_t^{k+1}}{k+1}
= 2\sum_{k=0}^\infty\frac{x_t^{2k+1}}{2k+1} 
= 2\sum_{k=0}^\infty\frac{1}{2k+1}\left(\frac{t-1}{t+1}\right)^{2k+1}\text.
\end{align*}
(2) Sei zunächst 
$\sigma:\!\  ]1, \infty[ \rightarrow \R, t \mapsto \frac{\sqrt t}{\log(t)}$.\\
Dann ist $\sigma$ offensichtlich stetig differenzierbar und es gilt für alle
$t \in\!\  ]1, \infty[$:
\[\sigma'(t)=\frac {\frac 1{2\sqrt t}\log(t)-\sqrt t\frac 1 t}{(\log(t))^2}
= \frac{\frac 1 {2\sqrt t}(\log(t) - 2)}{(\log(t))^2}\text.\]
Daher hat $\sigma'$ genau eine Nullstelle in $\exp(2)$, $\sigma$ ist auf 
$]1, \exp(2)[$ streng mototon
fallend, auf $]\exp(2), \infty[$ streng monoton wachsend und hat somit ein
Minimum in $\exp(2)$, sodass also
$\sigma(z)\ge\sigma(\exp(2))=\frac{\exp(1)}2$ für alle $z\in\!\  ]1,\infty[$
gilt.\\
Ferner gilt: $\exp(1)\ge 1 + 1 + \frac 1 2 + \frac 1 6 = \frac 8 3$ und somit
$\sigma(t) \ge \frac 1 2 \cdot\frac 8 3 = \frac 4 3$.\\
Somit gilt (die folgende Aussage ist für $x=1$ offensichtlich):
\begin{align}
\forall\,x\in [1,\infty[: \log(x) \le \frac 3 4 \sqrt x \text.
\label{nleq:2}
\end{align}
Sei weiter
$\alpha:\R_{>0} \rightarrow \R, s \mapsto \frac{s^{\frac 3 4}}{s+1}$.\\
Auch $\alpha$ ist stetig differenzierbar und es gilt für alle $s\in\R_{>0}$
\[\alpha'(s) = \frac{\frac 3 4 s^{-\frac 1 4}(s+1) - s^{\frac 3 4}}{(s+1)^2}
= \frac{\frac 34(s+1)-s}{s^{\frac 14}(s+1)^2}
=\frac{3-s}{4s^{\frac 1 4}(s+1)^2}\text.\]
Daher hat $\alpha'$ genau eine Nullstelle (in 3), $\alpha$ ist auf $]0, 3[$
streng monoton wachsend, auf $]3, \infty[$
streng monoton fallend.\\
Also gilt:
\begin{align}
 \alpha \text{ hat ein globales Maximum in }3\text.\label{nleq:3}
\end{align}
Ferner ist wegen
$3^5=243< \frac{24343800625}{100000000}=\left(\frac{395}{100}\right)^4$
auch 
\begin{align}3^{\frac 5 4} < \frac {395}{100}
\label{nleq:4}
\end{align}
Sei schließlich $\beta: \R_{>0} \rightarrow \R,
r \mapsto \frac 1 {r+1}\sqrt{\frac 3 4 r^2 + 2\pi r}$.\\
Dann ist $\beta$ stetig differenzierbar und es gilt für alle
$r \in \R_{>0}$
\begin{align*}
\beta'(r)&=\frac{\frac{\frac 64r+2\pi}{2\sqrt{\frac 34r^2+2\pi r}}(r+1)
 - \sqrt{\frac 3 4 r^2 + 2 \pi r}}
 {(r+1)^2}
 = \frac{\frac{\left(\frac 3 4 r + \pi\right)(r+1)- \frac 3 4 r^2 - 2\pi r}
              {\sqrt{\frac 3 4 r^2 + 2 \pi r}}}{(r+1)^2}\\
 &= \frac{\frac 3 4 r^2 + \pi r + \frac 3 4 r + \pi - \frac 3 4 r^2 - 2 \pi r}
 {(r+1)^2\sqrt{\frac 3 4 r^2 + 2 \pi r}}
= \frac{\left(\frac 34-\pi\right)r+\pi}{(r+1)^2\sqrt{\frac 34r^2+2\pi r}}\text.
\end{align*}
Sei $r_0 \coloneqq  \frac{\pi}{\pi -  \frac 3 4}$. Dann ist $\beta$ auf $]0,r_0[$
streng monoton wachsend, auf $]r_0, \infty[$ streng monoton fallend
und $\beta'$ hat genau eine Nullstelle (nämlich $r_0$).\\
Damit gilt:
\begin{align}
\beta \text{ hat ein globales Maximum in } r_0\text. \label{nleq:5}
\end{align}
Zuletzt gilt nach Archimedes:
\begin{align}\frac {223}{71} < \pi < \frac{22}7\text.
\label{nleq:6}
\end{align}
Damit ist auch wegen $6248\cdot 100 =624800<627396=4753\cdot 132$
\begin{align}\frac{6244}{4757} &=\frac{28\cdot 223}{67\cdot 71}
= \frac{\frac{223}{71}}{\frac {22} 7 - \frac 3 4}\notag\\ 
&< \frac{\pi}{\pi - \frac 3 4} < \frac{\frac {22}7}{\frac{223}{71}-\frac 3 4}
 = \frac{\frac {22}7}{\frac{679}{284}}= \frac{6248}{4753} < 
 \frac{132}{100}\text. \label{nleq:7} \end{align}
Sei nun $t \in [1, \infty[$. Dann gilt:
\begin{align*}
h_1(t)& = \frac{\sqrt t}{t+1}
\bigg((\pi-1)\sqrt{\log(t)}+\sqrt{\log(t)+2\pi}\bigg)\\
&\overset{\eqref{nleq:2}}\le
\frac{\sqrt t}{t+1}
\left((\pi-1)\sqrt{ \frac 3 4 \sqrt t }+\sqrt{\frac 3 4 \sqrt t +2\pi}\right)\\
&= (\pi-1)\sqrt{\frac 3 4}\frac {t^{\frac 3 4}}{t+1}
 + \frac 1 {t+1}\sqrt{\frac 3 4 t \sqrt t +  2\pi t}\\
&\overset{t \ge 1}\le(\pi-1)\sqrt{\frac 3 4}\frac {t^{\frac 3 4}}{t+1} + 
\frac 1 {t+1}\sqrt{\frac 3 4 t^2 +  2\pi t} 
= (\pi-1)\sqrt{\frac 3 4}\alpha(t) + \beta(t)\\
&\overset{\substack{\eqref{nleq:3},\\\eqref{nleq:5}}}\le 
(\pi-1)\sqrt{\frac 3 4}\alpha(3)
 + \beta\left(\frac{\pi}{\pi -  \frac 3 4} \right)\\
 &= (\pi-1)\frac 1 2 \cdot \frac{3^{\frac 5 4}}4 + 
 \frac 1 {\frac{\pi}{\pi -  \frac 3 4} + 1}
 \sqrt{\frac 3 4 \left(\frac{\pi} {\pi -  \frac 3 4} \right)^2 
 + 2\pi\frac{\pi}{\pi -  \frac 3 4}}\\
 &\overset{\substack{\eqref{nleq:6},\\\eqref{nleq:7}}}< 
 \frac{15}{56}3^{\frac 5 4} +
 \frac 1 {\frac{6244}{4757}+1}
 \sqrt{\frac 3 4 \cdot
  \frac{17424}{10000}+\frac{44}{7}\cdot\frac{132}{100}}\\
 &\overset{\eqref{nleq:4}}< \frac{5925}{5600} + \frac{4757}{11001}\sqrt{ 
  \frac{13068}{10000}+\frac{5808}{700}}
  < \frac{5925}{5600} + \frac{4757}{11001}
  \sqrt{\frac{13068}{10000}+\frac{58080}{7000}} \\
  &<\frac{237}{224} + \frac{4757}{11001}
  \sqrt{\frac{13068}{10000}+\frac{58086}{7000}}
  =\frac{237}{224} + \frac{4757}{11001}
  \sqrt{\frac{13068}{10000}+\frac{8298}{1000}}\\
  &=\frac{237}{224} + \frac{4757}{11001}
  \sqrt{\frac{96048}{10000}}
  < \frac{237}{224} + \frac{4757}{11001}
  \sqrt{\frac{96100}{10000}}
  = \frac{237}{224} + \frac{4757}{11001} \cdot \frac{31}{10}\\
  &= \frac{237}{224} + \frac{147467}{110010}
  < \frac{237}{224} + \frac{148}{110}
  = \frac{237}{224} + \frac{74}{55}
  =\frac{29611}{12320} < \frac{300}{123}\\
  &= \sqrt{\frac{90000}{15129}}\sqrt{\frac{71}{446}}\sqrt{\frac{446}{71}}
  = \sqrt{\frac{6390000}{6747534}}\sqrt{2\cdot \frac{223}{71}}
  < \sqrt{2\cdot \frac{223}{71}} \overset{\eqref{nleq:6}}< \sqrt{2\pi}\text.
\end{align*}
(3) Offenbar ist $h_2$ auf $\R_{>0}$ stetig differenzierbar und es gilt
für alle $t \in \R_{>0}$:
\begin{align}
&\,h_2'(t)\notag\\
 =&\, \frac{\frac t {2\sqrt{t+1}} - \sqrt{t+1}}{t^2}
\left((\pi-1)\sqrt{\log(1+t)} + \sqrt{\log(1+t)+2\pi} \right)\notag\\
&+ \frac{\sqrt{t+1}}t
\left(\frac{\pi-1}{2(t+1)\sqrt{\log(t+1)}} 
+ \frac 1 {2(t+1)\sqrt{\log(1+t) + 2\pi}} \right)\notag\\
=& \frac{-t-2}{2t^2\sqrt{t+1}}
  \left((\pi-1)\sqrt{\log(1+t)} + \sqrt{\log(1+t)+2\pi} \right)\notag\\
 & + \frac 1 {2t\sqrt{t+1}}\left(\frac{\pi-1}{\sqrt{\log(t+1)}} 
+ \frac 1 {\sqrt{\log(1+t) + 2\pi}}  \right)\notag\\
=&\,\scriptstyle{ \frac 1 {2t\sqrt{t+1}}
\left(- \frac {t+2}t\left((\pi-1)\sqrt{\log(1+t)} 
          + \sqrt{\log(1+t)+2\pi} \right)
          + \frac{\pi-1}{\sqrt{\log(t+1)}} 
+ \frac 1 {\sqrt{\log(1+t) + 2\pi}}  \right)}\text.
\label{nleq:8}
\end{align}
Nach (1) gilt für alle $x \in \R_{>0}: \log(x+1)\ge 2\frac{x}{x+2}$.\\
Sei $t \in \R_{>0}$. Dann gilt:
\[ \sqrt{\log(1+t)}\sqrt{\log(1+t)}=\log(1+t) \ge 2 \frac t{t+2} > 
\frac t {t+2}\text.\]
Also auch
\begin{align}
 \frac{1}{\sqrt{\log(1+t)}} < \frac{t+2}t\sqrt{\log(1+t)}\text.
 \label{nleq:9}
 \end{align}
Ferner gilt für alle $y \in [1, \infty[: -y + \frac 1 y \le 0$.\\
Es ist $y_t\coloneqq \sqrt{\log(1+t) + 2\pi}\ge 1$, also:
\begin{align} -\frac{t+2}t y_t + \frac 1 {y_t}
< - y_t + \frac 1 {y_t} < 0 \text.
\label{nleq:10}\end{align}
Somit gilt:
\begin{align*}
&2t\sqrt{t+1}\cdot h_2'(t)\\
& \overset{\eqref{nleq:8}}\le 
\scriptstyle{\left(- \frac {t+2}t\left((\pi-1)\sqrt{\log(1+t)} 
          + \sqrt{\log(1+t)+2\pi} \right)
          + \frac{\pi-1}{\sqrt{\log(t+1)}} 
+ \frac 1 {\sqrt{\log(1+t) + 2\pi}}  \right)}\\
& =(\pi-1)\underbrace{\left(\frac 1 {\sqrt{\log(1+t)}} 
- \frac{t+2}t\sqrt{\log(1+t)} \right)}_{\overset{\eqref{nleq:9}}<0}
+ \underbrace{\frac 1 {y_t} - \frac{t+2}ty_t}
             _{\overset{\eqref{nleq:10}}<0} < 0\text.
\end{align*}
Wegen $t > 0$ ist $h_2$ damit monoton fallend.\\
Es gilt:
\begin{align*}
 \exp\left(\frac{202}{100} \right) 
&\ge \sum_{k=0}^6\frac 1 {k!}\left(\frac{202}{100} \right)^k\\
&=\scriptstyle{1+\frac{202}{100} + \frac{40804}{20000}
+\frac{8242408}{6000000} + \frac{1664966416}{2400000000}
+ \frac{336323216032}{1200000000000}
+\frac{67937289638464}{720000000000000}}\\
&>\scriptstyle{1+\frac{202}{100} + \frac{204}{100}
+\frac{8242404}{6000000} + \frac{1664966400}{2400000000}
+ \frac{336323215920}{1200000000000}
+\frac{67937289638400}{720000000000000}}\\
&= \scriptstyle{1+\frac{202}{100} + \frac{204}{100}
+\frac{1373734}{1000000} + \frac{69373600}{100000000}
+ \frac{2802693466}{10000000000}
+\frac{94357346720}{1000000000000}}\\
&> 1+\frac{202}{100} + \frac{204}{100}
+\frac{1373}{1000} + \frac{693}{1000}
+ \frac{280}{1000}
+\frac{94}{1000} = \frac{7500}{1000}= \frac{15}2\text.
\end{align*}
Somit ist also 
\begin{align}
\log\left(\frac{15}{2}\right) < \frac{202}{100}\text.\label{nleq:11}
\end{align}
Weiter ist wegen $\frac {15} 2 < \frac{7502121}{1000000} =
\left(\frac{2739}{1000}\right)^2$ auch
\begin{align}
\sqrt{\frac{15}2}< \frac{2739}{1000}\text. \label{nleq:12}
\end{align}
Ebenso ist wegen $\frac{202}{100} < \frac{2022084}{1000000} =
\left(\frac{1422}{1000} \right)^2$ auch
$\sqrt{\frac{202}{100}}<\frac{1422}{1000}$
und analog ist auch wegen
$8305924\cdot 700 = 5814146800 > 5814000000 = 5814\cdot 1000000$
zunächst $\frac{5814}{700} < \frac{8305924}{1000000}=\left(\frac{2882}{1000}
\right)^2$ und damit
\begin{align}
\sqrt{\frac{5814}{700}}< \frac{2882}{1000}\text.\label{nleq:13}
\end{align}
Somit gilt wegen des monotonen Fallens von $h_2$ für alle 
$s \in [\frac{13}2, \infty[$:
\begin{align*}
h_2(s)&\le h_2\left(\frac{13}2\right) = \frac{\sqrt{\frac{15}2}}{\frac{13}{2}}
\left((\pi-1)\sqrt{\log\left(\frac{15}2 \right)} 
 + \sqrt{\log\left(\frac{15}2 \right) + 2\pi} \right)\\
 &\overset{\eqref{nleq:6}}<\frac{\sqrt{\frac{15}2}}{\frac{13}{2}}
\left(\frac{15}7\sqrt{\log\left(\frac{15}2 \right)} 
 + \sqrt{\log\left(\frac{15}2 \right) + \frac{44}7} \right)\\
& \overset{\substack{\eqref{nleq:11},\\\eqref{nleq:12}}}<\frac{5478}{13000}
 \left(\frac{15}7\sqrt{\frac{202}{100}} 
 + \sqrt{\frac{202}{100} + \frac{44}7} \right)\\
 &= \frac{5478}{13000}
 \left(\frac{15}7\sqrt{\frac{202}{100}}
 + \sqrt{\frac{5814}{700}}\right)
 \overset{\eqref{nleq:13}}< \frac{5478}{13000}
 \left(\frac{15}7 \cdot \frac{1422}{1000}
 + \frac{2882}{1000}\right)\\&= \frac{5478}{13000}\cdot\frac{41504}{7000}
 = \frac{227358912}{91000000}<\frac 5 2 
 =\sqrt{\frac{25}4 \cdot \frac{71}{446}}\sqrt{2 \frac{223}{71}}\\
 &= \sqrt{\frac{1775}{1784}}\sqrt{2 \frac{223}{71}}\overset{\eqref{nleq:6}}<
  \sqrt{2\pi}
 \text.
\end{align*}
(4) Seien $\eta \in\!\  ]0, 1[$ und $f_\eta:\!\  ]1, \infty[ \rightarrow \R,
t \mapsto \log(2t)-\eta t$.\\
Offenbar ist $f_\eta$ stetig differenzierbar und es gilt
$f_\eta'=\frac 1\cdot- \eta$.\\
Daher ist $\forall\, t \in [\frac 1 \eta, \infty[: f_\eta'(t) \le 0$.\\
Somit ist $\restrict{f_\eta}{[\nicefrac 1 \eta, \infty[}$ monoton fallend.\\
Wegen
\[4\log\left(\frac 2 \eta \right) = \log\left(\frac{16}{\eta^4} \right)
> \log(16) > 1\]
und damit $\frac 4 \eta \log\left(\frac 2 \eta\right) > \frac 1 \eta$ genügt es
also zu zeigen, dass
$f_\eta\left(\frac 4\eta\log\left(\frac 2 \eta\right)\right)<0$
gilt.\\
Es gilt allerdings $\forall\, x \in \R: \exp(x) \ge 1+x > x$, also
$\forall\, x \in \R: \log(x) < x$ und damit:
\begin{align*}
f_\eta\left(\frac 4 \eta \log\left(\frac 2 \eta\right) \right) & =
\log\left(\frac 8 \eta \log\left(\frac 2 \eta\right) \right)
 -4\log\left(\frac 2 \eta\right)\\
&= \log\left(\frac 8 \eta \log\left(\frac 2 \eta\right) \right)
  -\log\left(\frac{16}{\eta^4} \right)
  = \log\left(\frac{\frac{8 \log\left(\frac 2 \eta\right)}{\eta}}
    {\frac{16}{\eta^4}} \right)\\
&= \log\left(\frac{8\eta^4\log\left(\frac 2 \eta \right)}{16\eta} \right)
= \log\left(\frac 1 2 \eta^3 \log\left(\frac{2}{\eta} \right) \right)
< \log\left(\frac 1 2 \eta ^3 \frac 2 \eta \right)\\
& = \log(\eta^2) < 0\text.
\end{align*}
Somit gilt für
$t \in \bigg[\frac 4 \eta \log\left(\frac 2 \eta\right),\infty  \bigg[$ mit
$f_\eta(t)< 0$ auch $\log(2t)< \eta t$.
\end{proof}